\documentclass[12pt,letterpaper]{article}

% ---------- packages ----------
\usepackage[margin=1in,headheight=14.5pt]{geometry}
\usepackage{amsmath,amssymb}
\usepackage{booktabs}
\usepackage{array}
\usepackage{tabularx}
\usepackage{hyperref}
\usepackage{xcolor}
\definecolor{linkblue}{RGB}{70,90,160}
\usepackage{graphicx}
\usepackage{enumitem}
\usepackage{microtype}
\usepackage{parskip}
\usepackage{titlesec}
\usepackage{fancyhdr}
\usepackage{url}
\usepackage{cite}
\usepackage{caption}

% ---------- hyperref setup ----------
\hypersetup{
  colorlinks=true,
  linkcolor=black,
  citecolor=black,
  urlcolor=linkblue,
  pdftitle={COINjecture -- A Mathematics-Backed Peer-to-Peer Network},
}

% ---------- section formatting ----------
\titleformat{\section}{\large\bfseries}{}{0em}{}[\titlerule]
\titleformat{\subsection}{\normalsize\bfseries}{}{0em}{}
\titlespacing*{\section}{0pt}{1.5ex plus .2ex}{0.8ex plus .1ex}
\titlespacing*{\subsection}{0pt}{1ex plus .2ex}{0.4ex plus .1ex}

% ---------- header / footer ----------
\pagestyle{fancy}
\fancyhf{}
\renewcommand{\headrulewidth}{0.4pt}
\fancyhead[L]{\small COINjecture Network}
\fancyhead[R]{\small Whitepaper v2.6}
\fancyfoot[C]{\thepage}

% ---------- table column types ----------
\newcolumntype{L}[1]{>{\raggedright\arraybackslash}p{#1}}
\newcolumntype{C}[1]{>{\centering\arraybackslash}p{#1}}

% ---------- monospace url style ----------
\urlstyle{tt}

% ============================================================
\begin{document}

% ---------- title page ----------
\begin{titlepage}
  \centering
  \vspace*{2cm}

  {\LARGE\bfseries COINjecture Network}\\[0.4em]
  {\large A Mathematics-Backed Peer-to-Peer Network}\\[0.3em]
  {\normalsize \href{https://COINjecture.com}{COINjecture.com}}

  \vspace{1.5cm}
  {\normalsize Whitepaper\\Version 2.6, June 5, 2026}

  \vspace{2cm}

  % ---- Abstract ----
  \begin{minipage}{0.88\textwidth}
    \small\textbf{Abstract:}
    COINjecture is a purely peer-to-peer network that harnesses the solve-verify asymmetry of NP
    problems to replace the traditional proof-of-work hashing mechanism (e.g., SHA256) needed to
    outpace attackers and secure a robust network of digital payments. Chain validity derives from
    solution correctness and cumulative mathematical work, ensuring robust consensus without
    centralized validators. Difficulty is calibrated to instance hardness, not just solution validity,
    ensuring trivial solutions yield negligible rewards. The record formed in the process of solving
    and verifying these problems in the complexity class NP creates a computational marketplace
    where miners earn rewards for solving computationally hard problems while enabling secure,
    decentralized transactions.
  \end{minipage}

  \vspace{1.5cm}
  \begin{minipage}{0.88\textwidth}
    \small
    \textbf{Code Base:}\\
    Active:~\href{https://github.com/COINjecture-Network/COINjecture2.0}%
      {https://github.com/COINjecture-Network/COINjecture2.0}\\
    Archive:~\href{https://github.com/coinjecture-network/coinjecture}%
      {https://github.com/coinjecture-network/coinjecture}\\
    Archive:~\href{https://gitlab.com/beanapologist/coinjecture}%
      {https://gitlab.com/beanapologist/coinjecture}\\
    Archive:~\href{https://codeberg.org/beanapologist/COINjecture}%
      {https://codeberg.org/beanapologist/COINjecture}
  \end{minipage}

  \vfill
\end{titlepage}

% ============================================================
\section*{Introduction}

Traditional proof-of-work blockchain systems, like Bitcoin, rely on the Secure Hash Algorithm 256-bit
(SHA-256) as the mechanism to verify transactions, link blocks, and maintain network integrity and chain
immutability. The SHA-256 algorithm is applied to hash a block header repeatedly until one is found with
a specific prefix of leading zeros.~\cite{nakamoto} During this process a random number called a nonce
is changed over and over with each attempt in a brute-force effort to find the solution. This hash
calculation provides no scientific or mathematical value outside of network security and the computational
resources required to perform this brute-force search resulted in the consumption of over 150~TWh in
2024, according to the Cambridge Bitcoin Electricity Consumption Index.~\cite{cbeci} This computation
provides network security but no additional utility.

What is needed is a cryptographically secure electronic payment system that generates utility beyond
security. Building on Bitcoin's foundation of decentralized consensus through proof-of-work,
COINjecture proposes to replace brute-force hash search with NP problem solving as the core
proof-of-work mechanism that records verifiable, immutable computational data on-chain. Consensus is
derived from the mathematical correctness of submitted blocks and the overall work performed by the
chain. The computational work that fuels chain growth and provides network consensus also allows for
the creation of a computational marketplace in which miners are rewarded for solving computationally
hard problems submitted by users or generated by the network state.

This Proof of Useful Work (PoUW) model is theoretically grounded and has prior art in deployed
chains~\cite{ball2017,primecoin}. It has recently been extended to include NP problem optimization
specifically~\cite{oliver2017}. This peer-to-peer network directly rewards participants proportionally
based on their mathematical contributions with no fixed fee schedules or required transaction fees. Where
competitors utilize consensus tied to a fixed problem, COINjecture is designed to be extensible across a
registry of problem types and current testnet operations utilize subset sum, traveling salesman
problem (TSP), and Boolean Satisfiability Problems (SAT). As participation and demand grow,
additional NP problems will be added to the registry. There are theoretically thousands of problems
spanning areas including logistics, cryptography, scheduling and optimization, and biology, and
hundreds that have been mathematically proven to be part of the NP family.

% ============================================================
\section*{Transactions}

COINjecture follows the standard UTXO model in which we define an electronic coin as a chain of
digital signatures. This transaction format follows Bitcoin's UTXO model without modification beyond
the addition of a required proof bundle field in the block header. Each owner transfers the coin by
digitally signing a hash of the previous transaction and the public key of the next owner. In addition, all
blocks in our network must contain not just transactions, but also proofs of computational work (e.g., NP
problem solutions) that secure consensus and provide economic utility. This data alongside transactions is
broadcast to all nodes in the peer network.

% ============================================================
\section*{Timestamp Server}

Like Bitcoin, we employ a timestamp server that proves that the data must have existed at the time of
announcement. Each timestamp includes the previous timestamp in its hash as well as the computational
proofs (problem + solution) of the block, forming a chain. The proof bundle includes both a commitment
timestamp ($T_1$) and reveal timestamp ($T_2$) determined by the salt-commit-mine-reveal protocol
outlined in the subsequent section.

% ============================================================
\section*{Proof-of-Work}

To implement a distributed timestamp server on a peer-to-peer basis, we use a proof-of-work system
based on NP problems rather than hash collision. The defining feature of NP problems is that a candidate
solution can be verified as correct or incorrect, but not necessarily solved, in polynomial time---hence the
name \emph{nondeterministic polynomial time}---so that given a potential answer to the problem, you can
check if it is correct quickly, even if finding the answer in the first place is difficult or time-consuming.
For instance, a sudoku puzzle may take several minutes to solve, but can be checked for correctness
quickly by ensuring no row or column contains two occurrences of the same number.

\subsection*{a.\ Commitment Protocol}

To prevent grinding or generating many problems to find easy instances, we employ a succinct
salt-commit-mine-reveal protocol:

\begin{enumerate}[leftmargin=2em,itemsep=2pt]
  \item Miner commits to the problem: $\mathit{commitment} = H(\mathit{problem\_params} \,\|\, \mathit{salt})$
  \item Miner solves problem
  \item Miner computes: $\mathit{reveal} = H(\mathit{problem\_params} \,\|\, \mathit{salt} \,\|\, H(\mathit{solution}))$
  \item Miner publishes solution bundle
  \item If validated, new epoch salt deterministically derived from accepted block hash
\end{enumerate}

The epoch salt is deterministically derived from the parent block hash to prevent pre-mining activities.
Commitment cryptographically binds to solution via $H(\mathit{solution})$ while the hash reveals nothing
about the actual solution structure. As the system encourages multiple miners to compete on the same
problem instance, any miner who delays their commitment or artificially inflates their solve time risks
being outpaced by a competitor. This competitive pressure enforces honest timing without requiring a
trusted clock. Even if solve time is inflated, reward is bounded proportionally by the emission formula
(see \emph{Fee Schedules}).

\subsection*{b.\ Problem Types}

The commit-reveal protocol is designed to be extensible and allows the network to use different
NP-complete, co-NP-complete, and NP-hard problem types (any problem whose solution admits a short,
polynomial time-checkable certificate). Each problem type is included in a registry based on their
unique computational characteristics. These characteristics are then used to calculate the
difficulty and work score of the solution set. See Appendix~A for a complete problem registry and list of
proposed future additions.

\subsection*{c.\ Work Score Calculation}

Unlike traditional hash-based proof-of-work, where \texttt{work} is primarily the number of hashes attempted
(inferred from the difficulty target), our work score is a measurement of computational complexity that
takes into account solve and verification time, the correctness, and optimality of the solution. This work
score is not arbitrary and is constrained by network competition and verification of solutions.

The raw work score is anchored to the time asymmetry between solve time and verify time---the
fundamental characteristic that all NP problems share:
\[
  \mathit{raw\_work} = \log_2\!\left(\frac{\mathit{solve\_time}}{\mathit{verify\_time}}\right).
\]
The $\mathit{solve\_time}$ is measured as the block-timestamp difference between commitment inclusion
and solution reveal, anchored to the timestamp server's consensus chain. $\mathit{Verify\_time}$ is the
network-measured, instance-fixed denominator. Thus, the asymmetry is anchored to a real property of the
solution, not just wall-clock.

Using $\log_2$, we can measure computational difficulty in a unit comparable across problem types
without knowing anything about problem structure (\textit{problem\_params}). The complete work score
is time asymmetry with quality corrections as multiplicative adjustments:
\[
  \mathit{work\_score} = \log_2\!\left(\frac{\mathit{solve\_time}}{\mathit{verify\_time}}\right)
                         \times \mathit{quality\_score},
\]
where $\mathit{quality\_score} \in [0,1]$ (exact optimal solution $= 1$, valid but suboptimal $<1$).
Solution quality is established for each problem type in the problem registry and scored against the best
solution found in the current epoch. Both time asymmetry and solution quality are verifiable by the
network and cannot be gamed via self reporting.

\subsection*{d.\ Difficulty Adjustment}

The difficulty weight affects what happens \emph{before} mining (problem generation). The work score
measures what happened \emph{after} mining (solution quality). The difficulty target adjusts based on
the fixed window of $N = 20$ blocks to allow for stable analysis and to help support our optimal block
time of 10~seconds. The adjuster requires at least $N/2 = 10$ blocks of history before making its first
adjustment, and defers adjustment entirely when variance within the window is high ($\sigma > 0.8\mu$).
This value represents a tradeoff: larger $N$ produces smoother adjustment but slower response to
hashpower changes; smaller $N$ responds faster but risks oscillation.

At a 10-second block time, $N = 20$ corresponds to approximately 200~seconds of network
history---long enough for statistical stability, short enough to adapt within minutes. This difficulty
adjustment allows work scores to maintain stable block time and for decentralized participation in
the mining pool. The difficulty adjuster also prevents sandbagging by calibrating expectations to recent
network performance, so artificially inflated solve times produce scores that deviate from the target
without conferring a competitive advantage.

% ============================================================
\section*{Network}

The steps to run the network are as follows:

\begin{enumerate}[leftmargin=2em,itemsep=2pt]
  \item New transactions are broadcast to all nodes
  \item Miners select problems from epoch salts (user-submitted or consensus-generated)
  \item Miner binds to problem and creates commitment
  \item Miners solve problems
  \item When a valid solution and header are found, broadcast to all nodes
  \item All nodes validate the solution and accept the block if valid
  \item Nodes express acceptance by working on the next block using the salted hash of the accepted block
\end{enumerate}

Nodes always consider the chain with the highest cumulative work to be correct and work on extending
it. The cumulative work score is also the tiebreaker mechanism in the event that two nodes broadcast
different versions of the next block simultaneously. The asymmetry of NP problems (fast verification,
slow solving) means that nodes can quickly verify competing chains during a fork to determine which one
is valid and has the most ``weight'' according to the tiebreaker rule.

% ============================================================
\section*{Problem Pool}

The network maintains two problem sources: consensus problems and user problems. Consensus problems
are generated deterministically from network state for base rewards. User problems are submitted with
escrowed bounties. User problem work scores contribute to chain weight identically to consensus work.
Miners may allocate their compute resources freely between consensus problems (block rewards) and
marketplace problems (bounties).

% ============================================================
\section*{Incentive}

By convention, the first transaction in a block is a special transaction that starts a new coin owned by the
creator of the block. This adds an incentive for nodes to support the network and provides a way to
initially distribute coins into circulation, since there is no central authority to issue them.

Unlike Bitcoin, which uses the analogy of gold miners expending resources to add gold into circulation,
COINjecture incentive is derived through the measured performance of the nodes completing the work,
like a promotion based on the work you performed. In our case, it is empirically measured against the
actual performance of the chain and includes problem solution quality, solution speed, and the time it
takes to verify correctness. This quantifies ``good work'' and rewards it proportionally using the emission
formula:
\[
  \mathit{mint} = \left\lfloor \frac{w_{\mathrm{trunc}} \cdot S \cdot K}{\sqrt{W_{\mathrm{parent}}}} \right\rfloor,
\]
where $w_{\mathrm{trunc}}$ is the truncated work score of the current block, $K$ is the emission
multiplier (consensus constant $K=50$), $S$ is the fixed-point scale ($10^{12}$ atoms per display BEANS),
and $W_{\mathrm{parent}}$ is the sum of truncated work scores through the parent block (integer
$\lfloor\sqrt{W_{\mathrm{parent}}}\rfloor$ in the denominator on-chain). The emission multiplier, $K$, was 
selected to keep the average reward rate stable at 1~coins per block. This mint issuance is separate 
from fork-choice weight (cumulative $w_{\mathrm{trunc}}$) as reward does not affect security weight.

Each block's reward depends only on its own work and the parent's cumulative work score, with no future
dependence. The monotonicity created by proportional rewards based on work (more work yields more
reward) is proven in Lean~\cite{rewards_lean}. By design, you cannot inflate solve time without losing
blocks to faster miners. This tension between work score incentive and racing incentive (high time $=$
higher score vs.\ lower time $=$ block winner) further helps ensure honest participation in the system.

% ============================================================
\section*{Fee Schedules}

COINjecture does not utilise predetermined fee schedules. Instead, incentive emerges from the dynamics
of the decentralised network state itself. Consensus transactions carry no fees beyond an optional priority
tip because miners are already compensated by block rewards. User-submitted problems require escrowed
bounties, which serve as both payment for computational work and natural spam deterrence.
User-submitted problems also require a network fee upon submission. There are no ``wasted''
transactions because every submission either contributes to consensus or funds useful computation.

The emission formula creates nonzero asymptotic growth---as $W$ grows, per-block issuance falls. No
hard caps or limits, only the pressure of consensus reward decay. As the network matures, fee incentives
will shift from consensus problems as the primary driver to user-submitted bounties; however, rewards
for consensus transactions will always be compensated. As each block's mint is proportional to
$w_{\mathrm{trunc}} / \sqrt{W_{\mathrm{parent}}}$, even a successfully inflated work score would yield
diminishing marginal reward and an attacker optimising for score inflation gains an ever-smaller share of
a growing total.

These principles ensure the system adapts naturally to its own evolution without manual intervention and
that rewards scale dynamically. The reward structure may also help encourage nodes to stay honest and
drive difficulty higher. If a greedy attacker is able to recalculate all of the correct solutions and produce
the same quality of work as honest nodes, they would have to choose between using it to defraud people
by stealing back their payments, or using it to generate new coins. They ought to find it more profitable
to play by the rules---rules that favour them with more new coins than everyone else combined---than to
undermine the system and the validity of their own wealth.

Thus, the steady addition of useful computational work to the chain serves as both a security mechanism
and provides economic utility. Nodes are incentivised to perform and verify real mathematical work
quickly, creating value beyond network security.

% ============================================================
\section*{The Conjecture}

The preceding sections describe a designed network of competing miners whose collective behaviour is
shaped by the parameters established (e.g., solve times, difficulty adjustments, block production). The
conjecture is not that these quantities are fixed constants at genesis. It is that a network built from these
axioms and that measures and distributes rewards accordingly will empirically converge toward a
predicted equilibrium; both the COIN and conjecture in COINjecture. That prediction is falsifiable by
running the network and comparing measurements to the formulas in the Calculations section.

\subsection*{a.\ Design Axioms}

We establish two core design constraints or axioms for our system dynamics via symmetry and unit
normalisation. COINjecture tests whether a decentralised network built on these axioms naturally
converges to a mathematically predicted steady coherence state.

Let $\mu = x + iy$ represent the consensus eigenvalue where the unit-modulus constraint $|\mu|^2 = 1$
represents unitarity conservation.

\begin{enumerate}[leftmargin=2em,itemsep=2pt]
  \item Unit normalisation: $x^2 + y^2 = 1$
  \item Symmetry condition: $|x| = |y|$
\end{enumerate}

From symmetry, $x^2 = y^2$. Substituting into the unit constraint: $2x^2 = 1$, giving
$|x| = |y| = 1/\sqrt{2}$. Numeric checks of these axioms are verified at sample points in
\texttt{DesignAxioms.lean}~\cite{designaxioms_lean}. The constant $1/\sqrt{2}$ enters the protocol as
the difficulty target multiplier $\eta$, while the eigenvalue's phase is interpretive and does not feed
difficulty, reward, or fork choice.
The two axioms determine $|x| = |y| = 1/\sqrt{2}$ but leave four candidate points on the unit circle.
The whitepaper selects $\mu = (-1 + i)/\sqrt{2}$ from the additional constraint $\operatorname{Re} < 0$,
formalised as \texttt{mu\_unique} in \texttt{ComplexDecomposition.lean}~\cite{complexdecomp_lean}.
The modulus/argument decomposition used in the next subsection is the same file's anchor/coherence
factorisation.

\subsection*{b.\ Falsifiability}

The conjecture predicts the convergence $r = 1$ with transient deviations decaying as
$C(r^n) = \operatorname{sech}(n \cdot \log r)$, where the coherence function
$C(r) = 2r / (1 + r^2)$ attains its maximum $C(1) = 1$ at equilibrium. For our network, $r$ is the
dimensionless ratio between actual block solve time (the modulus) and the difficulty adjuster's optimal
target (the argument). This creates natural epistemic asymmetry: the numerator, block solve time, is
verified and fixed by the consensus chain once a block is accepted; the denominator, the difficulty
adjuster's optimal target, is self-referential and is recomputed from the network's own median.
Sech is the natural envelope for systems with reflective symmetry about equilibrium, where perturbations
in either direction are penalised identically---$C(r) = C(1/r)$. Each component plays an important role:
what the chain knows vs.\ what the chain expects. The conjecture is that the network's
measurement-and-reward dynamics drive this ratio toward 1.

The conjecture yields three testable predictions:

\begin{enumerate}[leftmargin=2em,itemsep=2pt]
  \item Perturbation recovery should follow the sech profile rather than exponential decay.
  \item The difficulty oscillation should exhibit a characteristic periodicity as the system converges to
        the lowest overhead state.
  \item The work score distribution should be symmetric about $\log(r) = 0$, reflecting $C(r) = C(1/r)$.
\end{enumerate}

Failure of any of these constitutes falsification. By building a network that measures itself and
distributes rewards accordingly (e.g., block solve times, difficulty scaling, optimality), we can test
whether the predicted coupling between resolved work and self-referential target sustains stable
operation.

% ============================================================
\section*{Reclaiming Disk Space}

Once the latest block in a chain is buried under enough blocks, solutions to problems can be discarded by
mining nodes to save disk space. To facilitate this without breaking the block's hash, solutions and
transactions are hashed in a Merkle tree. Block headers form the Merkle tree, with only the root included
in the block's hash. Old blocks can be compacted by stubbing off branches of the tree.

A block header with no transactions or solutions is about 1~KB. If we suppose that blocks are generated
at our optimal block time of 10~seconds---long enough for NP solving to be meaningful, short enough for
reasonable transaction finality, and enough overhead margin for verification plus propagation---header
storage is approximately 3.2~GB per year. With commodity servers commonly shipping with 1~TB or
more of storage as of 2025, this should not be a constraint even over decades of operation.

The network supports different node types with varying storage requirements:

\begin{itemize}[leftmargin=2em,itemsep=2pt]
  \item \textbf{Light Nodes:} Store only block headers and verify commitments. Can request full solutions
    from full nodes when needed. Suitable for mobile mining and everyday transactions.
  \item \textbf{Full Nodes:} Store complete blockchain and recent proof bundles (e.g., last $N$ epochs).
    Validate all transactions and solutions. Can prune old solutions after sufficient confirmations.
  \item \textbf{Archive Nodes:} Maintain complete historical record of all proofs. Enable academic
    verification, reproducibility, and serve as reference for the network.
\end{itemize}

Mining can occur on any node type. Light miners request work from full nodes, solve problems locally,
and submit solutions. This enables mining on resource-constrained devices while maintaining network
security through full node validation. For user-submitted problems, solutions are recorded on-chain but
full proof bundles can be pruned by non-archive nodes after the user retrieves their solution.

% ============================================================
\section*{Privacy}

Traditional privacy models rely on limiting access to information. COINjecture requires transparency of
measurements while preserving user privacy of methods.

\subsection*{a.\ Public Information (Required for Verification)}

The network must verify computational work, which requires:
\begin{itemize}[leftmargin=2em,itemsep=2pt]
  \item Problem parameters and solution
  \item Commitment timestamp ($T_1$) and reveal timestamp ($T_2$)
  \item Work score measurements ($\mathit{solve\_time}$, $\mathit{verify\_time}$, quality)
\end{itemize}
All nodes independently verify solution correctness and quality. This transparency is necessary for
trustless consensus.

\subsection*{b.\ Proprietary Information}

Computational methods remain proprietary while measurements remain verifiable. This preserves both
marketplace trust (public verification) and competitive advantage (private innovation). Transaction
privacy follows Bitcoin's model using pseudonymous addresses. The epoch salt, deterministically derived
from the parent block hash, prevents block linking to a common owner and ensures problems cannot be
pre-computed before commitment. This ensures that the network learns \textbf{WHAT} you solved
(problem and solution), \textbf{WHEN} you solved it ($T_1$ and $T_2$), and \textbf{HOW LONG} it
took ($\mathit{solve\_time} = T_2 - T_1$). The network never learns \textbf{HOW} you solved it
(algorithm) or \textbf{WHY} your implementation is efficient (proprietary optimisations).

% ============================================================
\section*{Calculations}

Unlike Bitcoin, which assumes each block adds the same amount of weight to the chain, each new block
in our chain has a different truncated work score. Consider an attacker controlling fraction $q$ of total
work-score production rate, with honest miners controlling $p = 1 - q$. Honest nodes verify before
accepting (rejecting any block which fails the deterministic checker) and follow the heaviest valid
tip---the chain with greatest cumulative truncated work $W = \sum w_{\mathrm{trunc}}$ from the common
ancestor. To replace the honest tip, an attacker must produce a valid alternate chain whose cumulative
$W$ exceeds the honest chain's.

The asymmetry of NP problems (exponential solve, polynomial verify) means nodes can quickly verify
competing chains during forks to determine cumulative work $W$. However, one challenge is quantifying
how much work an attacker must out-produce. The work score calculation utilises $\log_2$ to establish
a bit equivalent. Thus, cumulative $W$ is measured in bits and an attacker's deficit is a bit-deficit,
not a block-count deficit.

Where Bitcoin treats honest and attacker block production as competing Poisson processes whose event
rates are proportional to hash-power shares $p$ and $q$~\cite{nakamoto}, COINjecture treats them as
Poisson processes whose rates are proportional to work-score production shares, with $W$-bits replacing
block counts as the unit of ``length.'' An attacker who starts $z$ bit-confirmations behind must close that
bit-deficit before the honest chain extends further. When $p > q$, the probability of ever doing so drops
exponentially in $z$. The leading-order bound is an analogy of Bitcoin:
\[
  P(\text{attacker succeeds}) \approx \left(\frac{q}{p}\right)^{z}.
\]

% ---- Lean security bounds ----
The $P$ column is the simple bound $(q/p)^z$. Each additional honest confirmation multiplies the
attacker's failure probability by roughly $q/p$. \textbf{Lean proves the integer inequalities, not the
decimal cells.} For positive integers with $q < p$ and $z > 0$, Lean proves $q^z < p^z$---the statement
that an attacker's success probability is strictly below 1 at every confirmation depth~\cite{security_lean}.
Lean additionally proves that cumulative work strictly increases when a positive-work block is appended.
An attacker cannot fake $q$ as every block needs a verified solution, trivial problems score $\approx 0$,
and padding volume with self-funded bounties costs fees. Therefore, out-producing honest work is an
adversary's only path to rewriting history. That path closes exponentially as confirmation depth expands.

Each inequality below is Lean-verified~\cite{security_lean}. We note that NP solving, unlike SHA-256
hashing, is not memoryless---a miner mid-solve is closer to a solution than one who just started.
In practice, this means the Poisson bound is conservative; honest miners who have invested solve time
are harder to displace than the model assumes. As such, the probabilities are computed for display only
and will be more tightly bound in future work.

\begin{center}
\begin{tabular}{lllll}
\toprule
$q$ & $z$ & Bound & Lean & $P$ \\
\midrule
10\% & 5  & $(1/9)^5$    & $< 1$ & $\approx 1.7 \times 10^{-5}$  \\
10\% & 10 & $(1/9)^{10}$ & $< 1$ & $\approx 2.9 \times 10^{-10}$ \\
30\% & 10 & $(3/7)^{10}$ & $< 1$ & $\approx 2.1 \times 10^{-4}$  \\
49\% & 5  & $(49/51)^5$  & $< 1$ & $\approx 8.2 \times 10^{-1}$  \\
49\% & 20 & $(49/51)^{20}$ & $< 1$ & $\approx 4.5 \times 10^{-1}$ \\
\bottomrule
\end{tabular}
\end{center}

Therefore, COINjecture remains secure as long as no single solver out-produces everyone else across all
work types.

% ============================================================
\section*{Conclusion}

We have proposed a blockchain system that replaces traditional SHA-256 proof-of-work mining with NP
mathematical computations to create a marketplace for novel computational data. Nodes operate
independently with minimal coordination, yet every node can participate in the validation process
commensurate with resources. The mathematical work performed by the network serves as the consensus
mechanism and allows us to empirically measure the behaviours of a dynamical system in real time. This
creates the first blockchain where proof-of-work secures the network, provides economic utility, and
advances complexity research. This combination of incentive (COIN) and incomplete information
(conjecture) forms the basis for our peer-to-peer system.

% ============================================================
\section*{References}

\begin{thebibliography}{99}

\bibitem{nakamoto}
S.~Nakamoto, ``Bitcoin: A Peer-to-Peer Electronic Cash System,'' 2008.
\url{https://bitcoin.org/bitcoin.pdf}

\bibitem{cbeci}
Cambridge Centre for Alternative Finance, ``Cambridge Bitcoin Electricity Consumption Index (CBECI).''
\url{https://ccaf.io/cbnsi/cbeci}

\bibitem{ball2017}
M.~Ball, A.~Rosen, M.~Sabin, and P.~N. Vasudevan, ``Proofs of Useful Work,''
\textit{IACR Cryptology ePrint Archive}, Report 2017/203, 2017.
\url{https://eprint.iacr.org/2017/203}

\bibitem{primecoin}
S.~King, ``Primecoin: Cryptocurrency with Prime Number Proof-of-Work,'' 2013.
\url{https://primecoin.io/primecoin-paper.pdf}

\bibitem{oliver2017}
C.~G. Oliver, A.~Ricottone, and P.~Philippopoulos, ``Proposal for a Fully Decentralized Blockchain
and Proof-of-Work Algorithm for Solving NP-Complete Problems,'' arXiv:1708.09419, 2017.

\bibitem{rewards_lean}
COINjecture Network, ``Rewards.lean,'' COINjecture2.0 repository, 2026.
\url{https://github.com/COINjecture-Network/COINjecture2.0/blob/main/lean4/Coinjecture/Rewards.lean}

\bibitem{designaxioms_lean}
COINjecture Network, ``DesignAxioms.lean,'' COINjecture2.0 repository, 2026.
\url{https://github.com/COINjecture-Network/COINjecture2.0/blob/main/lean4/Coinjecture/DesignAxioms.lean}

\bibitem{complexdecomp_lean}
COINjecture Network, ``ComplexDecomposition.lean,'' COINjecture2.0 repository, 2026.
\url{https://github.com/COINjecture-Network/COINjecture2.0/blob/main/lean4/Coinjecture/ComplexDecomposition.lean}

\bibitem{security_lean}
COINjecture Network, ``SecurityCalculations.lean,'' COINjecture2.0 repository, 2026.
\url{https://github.com/COINjecture-Network/COINjecture2.0/blob/main/lean4/Coinjecture/SecurityCalculations.lean}

\bibitem{karp1972}
R.~M. Karp, ``Reducibility Among Combinatorial Problems,'' in \textit{Complexity of Computer
Computations}, Plenum Press, pp.~85--103, 1972.

\bibitem{horowitz1974}
E.~Horowitz and S.~Sahni, ``Computing Partitions with Applications to the Knapsack Problem,''
\textit{Journal of the ACM}, 21(2):277--292, 1974.

\bibitem{howgrave2010}
N.~Howgrave-Graham and A.~Joux, ``New Generic Algorithms for Hard Knapsacks,'' in
\textit{Advances in Cryptology -- EUROCRYPT 2010}, LNCS, vol.~6110, pp.~235--256, Springer, 2010.

\bibitem{held1962}
M.~Held and R.~M. Karp, ``A Dynamic Programming Approach to Sequencing Problems,''
\textit{Journal of the Society for Industrial and Applied Mathematics}, 10(1):196--210, 1962.

\bibitem{lin1973}
S.~Lin and B.~W. Kernighan, ``An Effective Heuristic Algorithm for the Traveling-Salesman Problem,''
\textit{Operations Research}, 21(2):498--516, 1973.

\bibitem{cook1971}
S.~A. Cook, ``The Complexity of Theorem-Proving Procedures,'' in \textit{Proceedings of the 3rd
Annual ACM Symposium on Theory of Computing}, pp.~151--158, 1971.

\bibitem{marques1999}
J.~P. Marques-Silva and K.~A. Sakallah, ``GRASP: A Search Algorithm for Propositional
Satisfiability,'' \textit{IEEE Transactions on Computers}, 48(5):506--521, 1999.

\bibitem{schoning1999}
U.~Sch\"{o}ning, ``A Probabilistic Algorithm for k-SAT and Constraint Satisfaction Problems,'' in
\textit{Proceedings of the 40th Annual IEEE Symposium on Foundations of Computer Science (FOCS)},
pp.~410--414, 1999.

\end{thebibliography}

% ============================================================
\appendix

\renewcommand{\thesection}{\Alph{section}}

\section{Problem Registry}

Our current testnet employs Subset Sum, Traveling Salesman Problem (TSP), and Boolean SAT. Each
problem type implements a \texttt{ProblemDescriptor} trait that declares its computational
characteristics. The difficulty adjuster and work score calculator operate generically over this
trait---adding a new problem type requires only a trait implementation and registration.

The registry's \texttt{base\_difficulty\_weight} informs the difficulty adjuster's cross-problem size
calibration but does not enter the work score formula, which remains problem-type agnostic.

\subsection*{a.\ Subset Sum}

Subset Sum is a classic NP-Complete problem with no known polynomial-time algorithm~\cite{karp1972}.
It asks if any subset of a given set of positive integers sums to a specific target value:
\[
  \text{Given } S = \{s_1, s_2, \ldots, s_n\} \text{ and target } T,
  \text{ find } S' \subseteq S \text{ such that } \sum_{s \in S'} s = T.
\]
Subset Sum serves as the canonical reference type---all size ratios and difficulty weights in the registry
are normalised relative to it. Verification is $O(n)$: sum the claimed subset and compare to $T$. Quality
is binary (correct $= 1.0$, incorrect $= 0.0$). Best known algorithms include Horowitz--Sahni
meet-in-the-middle $O(2^{n/2})$ and Howgrave-Graham--Joux $O(2^{0.337n})$~\cite{horowitz1974,howgrave2010}.

\subsection*{b.\ Traveling Salesman Problem (TSP)}

Given $n$ cities and a distance matrix, find the shortest Hamiltonian tour visiting each city exactly once
and returning to the start. NP-hard. Verification is $O(n)$: confirm the tour is a valid permutation, then
sum edge weights.

TSP is the primary vehicle for testing quality-weighted work scoring because it admits a quality
gradient---a valid tour that visits all cities is correct, but shorter tours score higher.
Quality $= 1 / (\mathit{tour\_length} + 1)$, scored relative to the best solution in the current epoch.
Best known exact algorithm is Held--Karp $O(2^n \cdot n^2)$; practical solvers use Lin--Kernighan
heuristics~\cite{held1962,lin1973}.

\subsection*{c.\ Boolean SAT}

Given a Boolean formula $\varphi$ in conjunctive normal form (CNF) with $n$ variables and $m$ clauses,
find a satisfying assignment. NP-complete (Cook--Levin theorem)~\cite{cook1971}. Verification is
$O(n \cdot m)$: substitute the assignment and confirm all clauses evaluate true.

COINjecture generates 3-CNF instances near the phase transition threshold ($\alpha \approx 4.27$ clauses
per variable), where instances are hardest. Quality is binary---partial satisfaction scores 0.0. Best known
algorithms include CDCL (Conflict-Driven Clause Learning) and Sch\"{o}ning's randomised
$O(2^{0.386n})$ for 3-SAT~\cite{marques1999,schoning1999}.

\subsection*{d.\ Registry Parameters (Currently Implemented)}

\begin{center}
\begin{tabular}{L{2.2cm}C{1.1cm}C{1.4cm}C{1.6cm}C{1.4cm}C{1.4cm}C{1.4cm}}
\toprule
Problem & Class & Scaling Exp. & Verify Cost & Size Ratio & Max Size & Quality \\
\midrule
Subset Sum  & NP-C & 0.80 & $O(n)$     & 1.00 & 180 & Binary   \\
SAT (3-CNF) & NP-C & 0.70 & $O(n{\cdot}m)$ & 1.00 & 320 & Binary   \\
TSP         & NP-H & 0.50 & $O(n)$     & 0.90 & 150 & Gradient \\
\bottomrule
\end{tabular}
\end{center}

\noindent\small NP-C = NP-Complete. NP-H = NP-Hard. Scale = scaling exponent.
Size R.\ = size relative to Subset Sum at equivalent difficulty.
Max = absolute safety ceiling, refined from observed testnet solve times.
\normalsize

\subsection*{e.\ Optimality (Current \& Future Problems)}

Decision problems (Subset Sum, SAT, Factorization) have binary quality: correct $= 1.0$, anything
else $= 0.0$. There is no gradient---partial solutions receive zero work score.

Optimization problems (TSP, Graph Coloring, SVP) have continuous quality. Valid solutions always score
above zero; better solutions score higher. Quality is scored relative to the best solution found in the
current epoch, making it a competitive measure. Invalid solutions (missing cities, adjacent same-colour
vertices, zero vector) score 0.0 regardless.

\begin{center}
\begin{tabular}{L{2.4cm}L{3.6cm}L{4.8cm}}
\toprule
Problem & Optimal Solution & Quality Formula \\
\midrule
Subset Sum    & Any $S'$ where $\sum s' = T$         & 1.0 if valid, else 0.0 \\
SAT           & Any satisfying assignment             & 1.0 if all clauses true, else 0.0 \\
TSP           & Shortest Hamiltonian tour             & $1/(\mathit{tour\_length}+1)$ vs.\ epoch-best \\
Graph Color.  & Fewest colours, no conflicts          & $1/\mathit{colors\_used}$ vs.\ epoch-best \\
Factorization & $p \times q = N$, non-trivial         & 1.0 if valid, else 0.0 \\
SVP           & Shortest non-zero lattice vector      & $1/(\|v\|+1)$ vs.\ epoch-best \\
\bottomrule
\end{tabular}
\end{center}

\subsection*{f.\ Future Problem Types}

Adding a new problem type requires implementing the \texttt{ProblemDescriptor} trait and calling
\texttt{register()}. No changes are required to the difficulty adjuster, work score calculator, or
consensus code. Other NP problems we intend to explore include:

\begin{center}
\begin{tabular}{L{2.2cm}C{1.1cm}C{1.4cm}C{1.6cm}C{1.4cm}C{1.4cm}C{1.4cm}}
\toprule
Problem & Class & Scaling Exp. & Verify Cost & Size Ratio & Max Size & Quality \\
\midrule
Graph Coloring & NP-C & 0.65 & $O(V{+}E)$ & 0.60 & 80  & Gradient \\
Factorization  & NP and co-NP & 0.33 & $O(M(\log N))$ & 0.50 & 200 & Binary   \\
SVP            & NP-H & 0.50 & $O(n^2)$     & 0.40 & 40  & Gradient \\
\bottomrule
\end{tabular}
\end{center}

% ============================================================
\section{Difficulty Adjustment}

The difficulty adjuster maintains stable block times by scaling problem size and is empirically tuned
based on observed solve times. All targets are derived from network state and mathematical
constants---no arbitrary parameters.

\begin{enumerate}[label=\alph*.,leftmargin=2em,itemsep=4pt]

\item \textbf{Target Times.} The optimal solve time is derived from the network's own median block time:
$\mathit{optimal} = \mathit{median\_block\_time} \times \eta$, where $\eta = 1/\sqrt{2}$. The
acceptable range is bounded by the golden ratio: $\mathit{min\_target} = \mathit{optimal} \times
\varphi^{-1}$, $\mathit{max\_target} = \mathit{optimal} \times \varphi$. During bootstrap (blocks
0--19), seeded defaults govern; after block~20, empirical data takes over.

\item \textbf{Adjustment Formula.} Given the moving average of the last $N$ block solve times, the
adjuster computes a dimensionless time ratio: $\mathit{time\_ratio} = \mathit{avg\_solve\_time} /
\mathit{optimal}$. Problem size scales as:
$\mathit{new\_size} = \mathit{current\_size} \times (1 / \mathit{time\_ratio})^{0.5}$,
clamped to $[0.4,\, 2.0]$ to prevent extreme jumps. The square root provides smooth convergence
without oscillation.

\item \textbf{High Variance Deferral.} If standard deviation exceeds 80\% of the mean ($\sigma >
0.8\mu$), the adjuster defers---it does not adjust difficulty until the window stabilises. This prevents
overreaction to transient spikes.

\item \textbf{Stall Recovery.} If average solve time exceeds $2 \times \mathit{max\_target}$, the
adjuster applies a stall penalty ($\mathit{size} \times 0.7$) and enters recovery mode, gradually
restoring size in increments of~1 once $\mathit{time\_ratio}$ falls below~1.2. Mining failures trigger a
softer penalty ($\mathit{size} \times 0.85$). These values have been empirically tuned during testnet
operation.

\end{enumerate}

% ============================================================
\section{User Submissions Integration}

The network maintains two problem sources: (1)~consensus problems generated deterministically from
network state for base rewards, and (2)~user-submitted problems with escrowed bounties. User
submissions create a computational marketplace where the mining network solves real problems for
payment.

Every submission requires an escrowed bounty, a minimum work score threshold, and an expiration
window. The bounty locks at submission and resolves in exactly one of three ways: released to a solver
on valid solution, refunded to the submitter on cancellation (open problems only), or refunded on
expiration. No partial claims. No intermediate access. Bounties finalise at confirmation depth, and depth
scales with bounty value ($W$-bits).

Users may select one of two submission modes.

\subsection*{a.\ Public Submissions}

Full problem specification broadcast at submission. Miners inspect, evaluate, and commit. Appropriate
when the problem carries no proprietary value.

\subsection*{b.\ Private Submissions}

Submitter publishes a cryptographic commitment, a well-formedness proof, and public parameters
including a complexity estimate. On the current testnet, this proof is a SHA-256 MAC placeholder.
A reveal phase is required before solutions are accepted---the submitter opens the commitment, the
network verifies the match, and solving proceeds as in public mode. A production zero-knowledge (ZK)
proof will ensure that the committed instance is valid and at the claimed difficulty is accurate without
revealing it. A formalized ZK proof is future work and is not yet deployed.

\subsection*{c.\ Aggregation Strategies}

\begin{center}
\begin{tabular}{L{2.2cm}L{3.8cm}L{4.8cm}}
\toprule
Strategy & Resolution & Use Case \\
\midrule
Any       & First valid solution                          & Decision problems (Subset Sum, SAT) \\
Best      & Highest quality in competition window         & Optimisation problems (TSP, Graph Coloring) \\
Multiple  & $N$ distinct valid solutions, bounty split    & Diverse solution sets \\
Statistical & $K$ solutions, aggregate statistics, bounty split & Research and distribution analysis \\
\bottomrule
\end{tabular}
\end{center}

% ============================================================
\section{Verify Lean Proofs Independently}

\begin{verbatim}
git clone https://github.com/COINjecture-Network/COINjecture2.0.git
cd COINjecture2.0/lean4
lake exe cache get
lake build
\end{verbatim}

The build succeeds with zero errors on Lean~4 v4.28.0 with Mathlib v4.28.0. The complete source is
available at:\\
\url{https://github.com/COINjecture-Network/COINjecture2.0/tree/main/lean4}

\end{document}
